\begin{landscape}
\begingroup
\fontsize{10}{13}\selectfont
\fancyhead[L]{\small NACHLASS. STÖRUNGEN DER PALLAS}
\renewcommand{\headrulewidth}{0pt}
\par\medskip
\SourceLeaf{563}
\pdfbookmark[2]{Nachlass VI [9]}{pallas-VI-fragment9}
\begin{center}[9.]\end{center}
\makeatletter
\providecommand{\VInineHistoricalNote}[1]{\begingroup
\renewcommand{\thefootnote}{*)}
\long\def\@makefntext##1{\noindent [*)\ ##1}
\footnote{#1}\endgroup}
\makeatother
\par\noindent\hspace*{\parindent}[Im Artikel 19 der Exposition ist bereits bemerkt worden, dass Pallas uns den Fall eines rationalen Verhältnisses der mittlern Bewegungen \(n\) und \(n'\) bietet; es verhalten sich nemlich die mittlere Bewegung von Pallas und Jupiter genau wie \(18:7\), so, dass sich dieses Verhältniss immer genau wiederherstellt\VInineHistoricalNote{Vgl. \textsc{Gauss} an \textsc{Bessel}, 5. Mai 1812, S. 421 dieses Bandes und Göttingische gelehrte Anzeigen vom 25. April 1812, Bd. VI, S. 349.]}. In diesem Artikel sollen die von dem Argument \(18\text{\mdseries\AstroSymbols ♃}-7\text{\mdseries\AstroSymbols ⚴}\) abhängigen Glieder für die practische Rechnung in säcularer Form gegeben werden, während die strenge Integration, durch welche die periodische Form wiederhergestellt wird, den Gegenstand des nächsten Artikels bilden wird.
\par\noindent\hspace*{\parindent}Aus den obigen numerischen Werthen für \(T,V,W\) erhält man nach Ausführung der Interpolation wie in Artikel 2—7 die von dem genannten Argument abhängigen Glieder
\[
\begin{aligned}
\text{in }\frac{(1-\cos i)\mathrm{d}\AscendingNode}{n\mathrm{d}t}:&\quad
0\overset{\prime\prime}{,}000\cos(18M'-7M)\\
&\qquad+0\overset{\prime\prime}{,}001\sin(18M'-7M)\\[.7ex]
\text{in }\frac{\mathrm{d}\tilde\omega}{n\mathrm{d}t}-\frac{(1-\cos i)\mathrm{d}\AscendingNode}{n\mathrm{d}t}:&\quad
+0\overset{\prime\prime}{,}095\cos(18M'-7M)\\
&\qquad-0\overset{\prime\prime}{,}055\sin(18M'-7M)\\[.7ex]
\text{in }\frac{\mathrm{d}\log\mathrm{hyp}\,n}{n\mathrm{d}t}=-\tfrac32\frac{\mathrm{d}\log\mathrm{hyp}\,a}{n\mathrm{d}t}:&\quad
+0\overset{\prime\prime}{,}012\cos(18M'-7M)\\
&\qquad-0\overset{\prime\prime}{,}027\sin(18M'-7M)\\[.7ex]
\text{in }\frac{\mathrm{d}\varepsilon}{n\mathrm{d}t}-\frac{(1-\cos i)\mathrm{d}\AscendingNode}{n\mathrm{d}t}:&\quad
-0\overset{\prime\prime}{,}039\cos(18M'-7M)\\
&\qquad-0\overset{\prime\prime}{,}021\sin(18M'-7M).
\end{aligned}
\]
\par\noindent\hspace*{\parindent}Indem man die beiden Glieder in \(\cos(18M'-7M)\) und \(\sin(18M'-7M)\) in Eines zusammenzieht und die mittlern Längen einführt durch die Werthe
\[
\begin{aligned}
M'&=\text{\mdseries\AstroSymbols ♃}-11^0 17'\quad5\overset{\prime\prime}{,}39\\
M&=\text{\mdseries\AstroSymbols ⚴}-121\quad8\quad54{,}50,
\end{aligned}
\]
erhält man das Glied
\[
\begin{aligned}
\text{in }\tilde\omega:&\quad-0\overset{\prime\prime}{,}1093\int\cos(18\text{\mdseries\AstroSymbols ♃}-7\text{\mdseries\AstroSymbols ⚴}+134^0 31'\,38'')\mathrm{d}\text{\mdseries\AstroSymbols ⚴}\\
\text{in }\log a:&\quad+0{,}415\int\cos(18\text{\mdseries\AstroSymbols ♃}-7\text{\mdseries\AstroSymbols ⚴}+170^0 56'\,57'')\mathrm{d}\text{\mdseries\AstroSymbols ⚴}\\
&\qquad\text{(in Einheiten der 7. Decimale)}\\
\text{in }\int n\mathrm{d}t:&\quad-0\overset{\prime\prime}{,}02955\iint\cos(18\text{\mdseries\AstroSymbols ♃}-7\text{\mdseries\AstroSymbols ⚴}+170^0 56'\,57'')\mathrm{d}\text{\mdseries\AstroSymbols ⚴}^{2}\\
\text{in }\varepsilon:&\quad-0\overset{\prime\prime}{,}0438\int\cos(18\text{\mdseries\AstroSymbols ♃}-7\text{\mdseries\AstroSymbols ⚴}+257^0 45'\,45'')\mathrm{d}\text{\mdseries\AstroSymbols ⚴};
\end{aligned}
\]
die Integration ergibt, für \(18\text{\mdseries\AstroSymbols ♃}-7\text{\mdseries\AstroSymbols ⚴}\) den constanten Werth \(125^0 36'\,49''\) gesetzt:
\begin{center}
\begin{tabular}{@{}cccll@{}}
Säcularänderung&von&\(\tilde\omega\)&\(:+0\overset{\prime\prime}{,}0255\)&(jährlich)\\
»&»&\(\log a\)&\(:+0{,}25\)&»\\
&&\multicolumn{3}{l}{(in Einheiten der 7. Decimale)}\\
»&»&\(\varepsilon\)&\(:-0\overset{\prime\prime}{,}0548\)&»\\
\multicolumn{3}{r}{Factor von \(tt\) in \(\int n\mathrm{d}t\)}&\(:-0\overset{\prime\prime}{,}0123\)&»
\end{tabular}
\end{center}
\par\noindent\hspace*{\parindent}Für die Neigung und die Länge des Knotens sind die entsprechenden Glieder sehr klein.]
\clearpage
\endgroup
\end{landscape}
